\begin{question}

  Suppose $\ket{i}$ and $\ket{j}$ are eigenkets of some Hermitian
  operator $A$. Under what condition can we conclude that $\ket{i} +
  \ket{j}$ is also an eigenket of $A$? Justify your answer.

\end{question}

\begin{answer}

  A is hermitian, therefore
  
  \begin{equation}
    \begin{split}
      &A\ket{a_i} = a_i\ket{a_i}\\
      &A\ket{a_j} = a_j\ket{a_j}\\
      &\bra{a_i}A = a_i\bra{a_i}\\
      &\bra{a_j}A = a_j\bra{a_j}\\
    \end{split}
  \end{equation}

  The characteristic relationship of Eigenkets must be respected, so
  we know that:

  \begin{equation}
    \begin{split}
      A(\ket{a_i} + \ket{a_j})
      &= A\ket{a_i} + A\ket{a_j}\\
      &= a_i\ket{a_i} + a_j\ket{a_j}\\
      &= a_k(\ket{a_i} + \ket{a_j})\\
    \end{split}
  \end{equation}

  The only way for this characteristic equation to be respected is
  if $a_i = a_j$ implying degeneracy.
  
\end{answer}
